\section{Discussion}
\label{sec:discussion}

The controlled comparisons answer two linked questions.  First, within each
architecture, does a downloaded SAR-domain initialization improve
label-efficient vessel detection relative to optical remote-sensing,
ImageNet, and random initialization?  Second, does the direction of that
within-track contrast persist across ViT and CNN backbones?  We interpret
Eqs.~\eqref{eq:transfer-gain} and~\eqref{eq:sar-optical-gap} in that order.
Cross-track agreement strengthens an architecture-generality claim, but a
difference between tracks cannot be attributed to architecture alone without
the qualifications below.

\ifdiagnosticdraft
Because the corrected H100 cohort and once-only evaluation are pending, this
draft does not select a winning initialization or claim a source-domain
advantage.  Both a consistent SAR advantage and a null or architecture-specific
pattern would answer the stated research question; the discussion will be
completed from the guarded finding macros rather than from V100 diagnostics.
\else
\FindingArchitectureGeneralization
\fi

\subsection{Limits of the comparison}
\label{sec:limitations}

\paragraph{One fixed seed.}
All 32 controlled cells use seed~0.  The curves are therefore point estimates,
not estimates of run-to-run variability.  We report no error bars,
confidence intervals, or statistical significance tests, and small
fraction-wise differences should not be overinterpreted.

\paragraph{In-region rather than geographic holdout.}
The splits are disjoint by scene identifier, but Sentinel-1 revisits create
geographic proximity across splits: 22 of 23 dev scenes fall within 32~km of
a training scene, and 42 of the 50 human-verified scenes fall within 50~km of
a training scene.  The reported metrics therefore characterize in-region
generalization under scene-level separation, not transfer to wholly unseen
geographies or acquisition regimes.

\paragraph{Uneven slice support.}
The test split contains 1,165 valid vessel targets but only two valid
near-shore vessel targets.  Near-shore test F1 is correspondingly sparse.
Dark-vessel labels are available only in the disjoint 50-scene
human-verified set, which is evaluated once after model and threshold
freezing.  Neither slice should be extrapolated beyond its observed support.

\paragraph{Source domain is not isolated from pretraining recipe.}
The ViT comparison pairs SatDINO, trained with a DINO objective on optical
fMoW imagery, with SARMAE, trained by masked autoencoding on SAR-1M.  Their
difference includes objective and corpus construction as well as modality.
The CNN BigEarthNet Sentinel-2/Sentinel-1 pair is the cleaner modality
contrast because model family, downstream source task, and dataset family are
matched more closely.  The generic controls also have different histories:
the ViT checkpoint uses supervised AugReg on ImageNet-1K, whereas the
ConvNeXt-V2 checkpoint uses FCMAE pretraining followed by supervised
ImageNet-1K fine-tuning.  We therefore make domain claims within architecture
and treat cross-track agreement as replication of a pattern, not a pooled
causal estimate of source modality.

\paragraph{Size matching is not compute matching.}
ViT-B/16 and ConvNeXt-V2-Base are approximately size matched, but their
pretraining data volume, optimization compute, token or feature geometry, and
training histories differ.  Their detector adapters are necessarily
family-specific even though both emit the same stride-4 output geometry.
Consequently, matched-role ViT--CNN gaps include architecture, adapter, and
pretraining-history effects.  Measured fine-tuning GPU-hours and parameter
counts make these differences visible but do not eliminate them.

\paragraph{Numerical and hardware scope.}
Every core cell uses one process per H100, micro-batch and effective batch 16,
Lightning \texttt{32-true}, and disabled CUDA-matmul and cuDNN TF32.  This
strict FP32 choice removes precision as a within-study variable and follows a
non-finite mixed-precision diagnostic, but it does not establish that FP32 is
universally preferable or that identical rankings would hold under another
numerical format or hardware generation.

\paragraph{External references are contextual.}
YOLO26 and LocateAnything are valuable operational reference points, but they
do not share the controlled backbones, pretraining roles, input handling, or
training protocol.  They remain outside the core curves and cannot be used to
infer the label savings or source-domain effect defined by this experiment.
Finally, the study covers one dataset, one point-detector family, four tested
label budgets, and one operating-seed realization; broader claims require
additional sensors, regions, detectors, and seeds.
